\subsection{上下界网络流}

\subsubsection{上下界可行流/最大流/最小流}

基于 \texttt{Dinic}（见 MaxFlow.tex，其中定义了 \texttt{i64} 与 \texttt{INF}）。\texttt{connect(u,v,l,r)}：添加边 $u\to v$，流量须落在区间 $[l,r]$ 内。\texttt{getFlow()} 判断是否存在满足所有边上下界的可行流（不带源汇即判断是否存在循环流）；\texttt{maxflow(s,t)}、\texttt{minflow(s,t)} 返回满足边上下界的 $s\to t$ 最大/最小流量，不存在可行流时返回 $-1$。该方法求得的可行流值非负，即需要存在一条流量值 $\ge 0$ 的可行流。\texttt{maxflow} 与 \texttt{minflow} 各在可行流判定后多做一次增广，总复杂度与两次 \texttt{Dinic} 同阶，$\mathcal{O}(n^2m)$。

\begin{lstlisting}
template<typename Flow>
struct BoundFlow {
    Flow G;
    std::vector<i64> w,f;
    int n,m = 0;
    BoundFlow (int _n) : n(_n),G(_n + 2),w(_n+1),f(1) {}
    void connect(int u,int v,i64 l,i64 r) {
        m++;
        G.connect(u,v,r-l);
        f.push_back(l);
        w[u] -= l;
        w[v] += l;
    }
    bool getFlow(int s = 0,int t = 0) {
        int S = n+1,T = n+2;
        for (int i = 1;i <= n;i++){
            if (w[i] > 0) {
                G.connect(S,i,abs(w[i]));
            } else if (w[i] < 0) {
                G.connect(i,T,abs(w[i]));
            }
        }
        int M = G.edge.size();
        if (s && t) G.connect(t,s,INF);
        G.work(S,T);
        for (int i = 2;i <= m*2;i += 2){
            f[i/2] += G.edge[i].c - G.edge[i].f;
        }
        for (int i = m*2+2;i < M;i += 2){
            if (G.edge[i].f != 0 || G.edge[i^1].f != G.edge[i^1].c) return 0;
        }
        return 1;
    }
    i64 maxflow(int s,int t){
        if (!getFlow(s,t)) return -1;
        int rid = G.head[t];
        i64 sum = G.edge[rid^1].f;
        G.edge[rid].f = G.edge[rid^1].f = 0;
        G.edge[rid].c = G.edge[rid^1].c = 0;
        sum += G.work(s,t);
        return sum;
    }
    i64 minflow(int s,int t){
        if (!getFlow(s,t)) return -1;
        int rid = G.head[t];
        i64 sum = G.edge[rid^1].f;
        G.edge[rid].f = G.edge[rid^1].f = 0;
        G.edge[rid].c = G.edge[rid^1].c = 0;
        sum -= G.work(t,s);
        return sum;
    }
};
\end{lstlisting}